\begin{lemma}[The terminal atom lies in the original support]
If $h$ is bad for $\mathsf{PowerRel}(r)$, then for every $x$ there are
$k,q,q'$ such that the two adjacent copied values are
$\mathsf{atom}(q)$ and $\mathsf{atom}(q')$ and
\[
q\in\mathsf{PowerSupport}(h(x)).
\]
\end{lemma}
\begin{proof}
Fix $x$.  Apply the termination statement
\noderef{PowerCopiedTerminates} with $n=0$ to obtain a depth $k$ such
that
\[
\mathsf{PowerCopiedTerminalAt}_r(h,x,0,k)
\]
holds.  By the definition of this terminal condition
\noderef{PowerCopiedTerminalAt}, unpacking its existential gives atoms
$q,q'$ with
\[
\mathsf{PowerCopiedLeft}_r(h,x,0,k)=\mathsf{atom}(q),
\qquad
\mathsf{PowerCopiedLeft}_r(h,x,0+1,k)=\mathsf{atom}(q').
\]
Thus the two adjacent copied values have the required atom forms.

Now apply the support containment
\noderef{PowerCopiedSupport} at row $0$ and depth $k$:
\[
\mathsf{PowerSupport}(\mathsf{PowerCopiedLeft}_r(h,x,0,k))
\subseteq
\mathsf{PowerSupport}(h(\mathsf{PowerShiftIter}(0,x))).
\]
The zeroth-iterate equation in \noderef{PowerShiftIter} says
$\mathsf{PowerShiftIter}(0,x)=x$.  Substituting the first atom equation
and this zeroth-iterate equation therefore gives
\[
\mathsf{PowerSupport}(\mathsf{atom}(q))
\subseteq \mathsf{PowerSupport}(h(x)).
\]
Finally, the atom clause of the support definition \noderef{PowerSupport}
identifies
\[
\mathsf{PowerSupport}(\mathsf{atom}(q))=\{q\},
\]
so $q\in\mathsf{PowerSupport}(\mathsf{atom}(q))$.  Applying the displayed
subset inclusion to this membership yields
$q\in\mathsf{PowerSupport}(h(x))$ by explicit transitivity of membership
through a subset.  Together with the two atom equations above, this is the
claimed existential.
\end{proof}
